\documentclass[aspectratio=169]{beamer}

% --- Theme and Style ---
\usetheme{Madrid}
\usecolortheme{whale}
\usefonttheme{professionalfonts}

% --- Packages ---
\usepackage[utf8]{inputenc}
\usepackage[T1]{fontenc}
\usepackage[indonesian]{babel}
\usepackage{amsmath, amsfonts, amssymb}
\usepackage{graphicx}
\usepackage{booktabs}
\usepackage{tikz}
\usepackage{caption}
\usepackage{subcaption}
\usepackage{algorithm}
\usepackage{algpseudocode}

% --- Metadata ---
\title[Integrasi EPG  dan   VQE]{Integrasi Exact Potential Game dan  Variational Quantum Eigensolver untuk Optimasi Portofolio Hibrida}
\subtitle{Tugas Akhir - Econophysics  dan   Quantum Finance}
\author[Nama Anda]{Nama Anda} % Silakan ganti dengan nama lengkap
\institute[Institusi]{Pembimbing: Nama Pembimbing \\ \medskip Departemen Fisika \\ Universitas ...}
\date{\today}

\graphicspath{{../Images/}}

% --- Custom Colors (Optional) ---
\definecolor{primaryblue}{RGB}{0, 31, 63}
\setbeamercolor{palette primary}{bg=primaryblue, fg=white}
\setbeamercolor{titlelike}{parent=palette primary}

% --- Logo ---
%\logo{\includegraphics[height=1cm]{Images/logo_universitas.png}}

\begin{document}

% Slide 1: Judul
\begin{frame}
    \titlepage
\end{frame}

% Daftar Isi
\begin{frame}{Outline}
    \tableofcontents
\end{frame}

% --- Konten Modular ---
\section{Pendahuluan}

% Slide 2: Latar Belakang  dan   Motivasi
\begin{frame}{Latar Belakang  dan   Motivasi}

    \begin{itemize}
        \item \textbf{Limitasi Model Klasik (Markowitz):}
        \begin{itemize}
            \item Asumsi distribusi normal dan  stasioneritas seringkali tidak terpenuhi.
            \item Kesulitan menangkap korelasi non-linear selama krisis pasar.
        \end{itemize}
        \item \textbf{Komputasi Kuantum:}
        \begin{itemize}
            \item Transformasi masalah optimasi ke dalam model \textit{Ising} / QUBO.
            \item Eksplorasi ruang solusi besar secara paralel pada perangkat NISQ.
        \end{itemize}
        \item \textbf{Game Theory (EPG):}
        \begin{itemize}
            \item Pemodelan interaksi strategis multi-agen di pasar finansial.
            \item Penyederhanaan optimasi melalui fungsi potensial global.
        \end{itemize}
    \end{itemize}
\end{frame}

% Slide 3: Rumusan Masalah  dan   Batasan
\begin{frame}{Rumusan Masalah  dan   Batasan}
    \begin{block}{Rumusan Masalah}
        \begin{enumerate}
            \item Bagaimana pengaruh formulasi EPG terhadap akurasi pemetaan Hamiltonian?
            \item Seberapa efektif strategi \textit{warm-start} (Nash) dalam mengatasi \textit{Barren Plateaus}?
            \item Bagaimana performa VQE dengan SPSA dan  \textit{Ansatz} EfficientSU(2)?
        \end{enumerate}
    \end{block}
    \begin{block}{Batasan Masalah}
        \begin{itemize}
            \item Portofolio: $N=2$ (BBCA, TLKM)  dan   $N=4$ aset.
            \item Framework: \textit{Pennylane} simulator.
            \item Algoritma Ekuilibrium: \textit{Sequential Best Response} (SBR).
        </itemize>
    \end{block}
\end{frame}
\section{Landasan Konsep Finansial}

% Slide 4: Bab 2.1 - Dinamika Pasar Finansial
\begin{frame}{Bab 2.1 - Dinamika Pasar Finansial}
    \begin{itemize}
        \item \textbf{Evolusi Uang:} Transformasi dari kebutuhan biologis (barter) $\rightarrow$ komoditas $\rightarrow$ uang \textit{fiat} $\rightarrow$ entitas data digital.
        \item \textbf{Inflasi:} Fenomena sistemik peningkatan entropi ekonomi; degradasi daya beli yang memaksa realokasi modal strategis.
        \item \textbf{Investasi:} Pengorbanan konsumsi jangka pendek untuk imbal hasil masa depan (\textit{Risk-Return Trade-off}).
        \item \textbf{Risk Aversion:} Parameter endogen yang dipengaruhi oleh Teori Prospek (respon asimetris terhadap kerugian).
    \end{itemize}
\end{frame}

% Slide 5: Bab 2.2 - Representasi Data Finansial
\begin{frame}{Bab 2.2 - Representasi Data Finansial}
    \begin{itemize}
        \item \textbf{Analisis Kualitatif (Fundamental):} Evaluasi kesehatan industri, laporan keuangan (\textit{Cashflow}), dan  sentimen pasar.
        \item \textbf{Analisis Teknikal (Visual):} Identifikasi pola memori sosial melalui \textit{Candlestick} dan  level \textit{Support/Resistance}.
        \item \textbf{Analisis Kuantitatif:}
        \begin{itemize}
            \item Transformasi \textit{Log-Return} untuk mencapai stasioneritas data.
            \item Normalisasi rentang nilai menggunakan fungsi aktivasi ($tanh$, $Sigmoid$).
        \end{itemize}
        \item \textbf{Model Markowitz:} Formulasi Lagrangian untuk optimasi \textit{Mean-Variance} portofolio.
    \end{itemize}
\end{frame}
\section{Teori Permainan  dan   Model Ising}

% Slide 6: Bab 2.3 - Teori Permainan  dan   EPG
\begin{frame}{Bab 2.3 - Teori Permainan  dan   EPG}
    \begin{itemize}
        \item \textbf{Formalisme Aksiomatik:} Logika strategi di bawah kondisi ketergantungan sistemik (Von Neumann  dan   Morgenstern).
        \item \textbf{Exact Potential Game (EPG):} Penyelarasan utilitas marginal agen dengan fungsi potensial global $\Phi$.
        \item \textbf{Ekuilibrium Nash:} Kondisi stabil di mana tidak ada insentif untuk deviasi strategis sepihak.
        \item \textbf{Markowitz sebagai EPG:} Pembuktian bahwa minimisasi risiko portofolio biner setara dengan maksimisasi utilitas EPG.
    \end{itemize}
\end{frame}

% Slide 7: Bab 2.4 - Hamiltonian Ising
\begin{frame}{Bab 2.4 - Hamiltonian Ising}
    \begin{itemize}
        \item \textbf{Sejarah Model Ising:} Investigasi mikroskopis feromagnetisme (Lenz  dan   Ising, 1920-an).
        \item \textbf{Representasi Aset:} Aset sebagai partikel \textit{spin} ($s_i \in \{-1, +1\}$) dalam kisi kristal terstruktur.
        \item \textbf{Parameter Fisik:}
        \begin{itemize}
            \item Me dan  Lokal ($h_i$): Representasi ekspektasi imbal hasil aset.
            \item Kopling ($J_{ij}$): Representasi korelasi interaksi risiko antar aset.
        \end{itemize}
        \item \textbf{Ground State:} Solusi optimal portofolio setara dengan konfigurasi energi terendah sistem.
    \end{itemize}
\end{frame}
\section{Komputer Kuantum}

% Slide 8: Bab 2.5 - Fondasi Komputer Kuantum
\begin{frame}{Bab 2.5 - Fondasi Komputer Kuantum}
    \begin{itemize}
        \item \textbf{Sistem Qubit:} Perluasan bit klasik menjadi status superposisi $|\psi\rangle = \alpha|0\rangle + \beta|1\rangle$ (Bola Bloch).
        \item \textbf{Evolusi Uniter:} Manipulasi status melalui gerbang rotasi ($R_x, R_y, R_z$) dan  \textit{entanglement} (CNOT).
        \item \textbf{Pengukuran (Born Rule):} Kolaps fungsi gelombang menjadi probabilitas deterministik $|\alpha|^2$ and $|\beta|^2$.
        \item \textbf{Parameterized Quantum Circuit (PQC):} Arsitektur sirkuit adaptif untuk optimasi variasional (Ansatz).
    \end{itemize}
\end{frame}
\section{Optimasi  dan   VQE}

% Slide 9: Bab 2.6 - Optimasi Parameter Klasik
\begin{frame}{Bab 2.6 - Optimasi Parameter Klasik}
    \begin{itemize}
        \item \textbf{Parameter Shift Rule:} Estimator gradien analitik eksak tanpa bias untuk sirkuit kuantum.
        \item \textbf{Gradient Descent:} Penurunan parameter menuju minimum global dengan \textit{learning rate} $\eta$.
        \item \textbf{Optimizer SPSA:}
        \begin{itemize}
            \item \textit{Simultaneous Perturbation Stochastic Approximation}.
            \item Reduksi kompleksitas dari $\mathcal{O}(p)$ menjadi $\mathcal{O}(1)$ (hanya 2 pengukuran per iterasi).
            \item Tangguh terhadap noise pada perangkat NISQ.
        \end{itemize}
    \end{itemize}
\end{frame}

% Slide 10: Bab 2.7 - Algoritma VQE
\begin{frame}{Bab 2.7 - Algoritma VQE}
    \begin{itemize}
        \item \textbf{Prinsip Variasional:} Teorema Rayleigh-Ritz untuk estimasi energi \textit{ground state} $E_0$.
        \item \textbf{Unitary Coupled Cluster (UCC):} Fondasi Ansatz kimia kuantum (eksponensial uniter).
        \item \textbf{EfficientSU2:} Arsitektur \textit{Hardware-Efficient} menggunakan lapisan rotasi dan  keterbelitan linear (CX).
        \item \textbf{Siklus Hibrida:} Distribusi beban komputasi antara QPU (evaluasi energi) dan  CPU (update parameter).
    \end{itemize}
\end{frame}

% Slide 11: Bab 2.8 - Validasi Barren Plateau
\begin{frame}{Bab 2.8 - Validasi Barren Plateau}
    \begin{itemize}
        \item \textbf{Barren Plateau:} Fenomena lenyapnya gradien secara eksponensial seiring bertambahnya jumlah qubit.
        \item \textbf{Penyebab:} Inisialisasi parameter acak pada sirkuit dengan \textit{entanglement} tinggi.
        \item \textbf{Entropi Von Neumann:} Metrik untuk mengukur derajat keacakan dan  keterbelitan state kuantum.
        \item \textbf{Solusi (Proposed):} Inisialisasi EPG/Nash untuk memandu parameter keluar dari zona tandus.
    \end{itemize}
\end{frame}
\section{Metodologi}

% Slide 12: Alur Metodologi Hibrida
\begin{frame}{Alur Metodologi Hibrida}
    \begin{center}
        \begin{tikzpicture}[node distance=2cm, auto, thick]
            \node [draw, rectangle] (data) {Data Harga};
            \node [draw, rectangle, right of=data, xshift=1.5cm] (epg) {EPG (Nash)};
            \node [draw, rectangle, right of=epg, xshift=1.5cm] (ham) {Hamiltonian};
            \node [draw, rectangle, below of=ham] (vqe) {VQE Optim.};
            \node [draw, rectangle, left of=vqe, xshift=-1.5cm] (res) {Optimal Portf.};
            
            \draw [->] (data) -- (epg);
            \draw [->] (epg) -- (ham);
            \draw [->] (ham) -- (vqe);
            \draw [->] (vqe) -- (res);
        \end{tikzpicture}
    \end{center}
    \begin{itemize}
        \item Integrasi Teori Permainan sebagai \textbf{Warm-Start} strategis untuk sirkuit kuantum.
    \end{itemize}
\end{frame}

% Slide 13: Formulasi Hamiltonian Ising Spesifik
\begin{frame}{Formulasi Hamiltonian Ising Spesifik}
    \begin{equation}
        H = \sum_i h_i \sigma_i^z + \sum_{i<j} J_{ij} \sigma_i^z \sigma_j^z + A(\text{Penalty})
    \end{equation}
    \begin{itemize}
        \item \textbf{Me dan  Lokal ($h_i$):} Menyerap bias imbal hasil and risiko individual.
        \item \textbf{Kopling ($J_{ij}$):} Menangkap dinamika korelasi (risiko bersama) antar aset.
        \item \textbf{Suku Penalti ($A$):} Menjamin kepatuhan terhadap batasan anggaran (budget constraint).
    \end{itemize}
\end{frame}

% Slide 14: Konfigurasi Sirkuit  dan   Optimizer
\begin{frame}{Konfigurasi Sirkuit  dan   Optimizer}
    \begin{itemize}
        \item \textbf{Framework:} PennyLane dengan simulator \texttt{default.qubit}.
        \item \textbf{Hiperparameter SPSA:}
        \begin{itemize}
            \item $a = 0.01, c = 0.1$ (Laju pembelajaran and perturbasi).
            \item Iterasi: 200 step per rebalancing.
        \end{itemize}
        \item \textbf{Ansatz:} \texttt{EfficientSU2} dengan kedalaman $L=2$.
    \end{itemize}
\end{frame}
\section{Hasil dan  Pembahasan}

% Slide 15: Strategi Warm-Start  dan   Nash Equilibrium
\begin{frame}{Strategi Warm-Start  dan   Nash Equilibrium}
    \begin{itemize}
        \item PSNE (\textit{Pure Strategy Nash Equilibrium}) diinjeksikan sebagai inisialisasi parameter $\theta_0$.
        \item \textbf{Hasil:} Memandu sirkuit melewati area \textit{Barren Plateaus} sejak iterasi awal.
        \item Mempercepat konvergensi energi menuju \textit{ground state}.
    \end{itemize}
\end{frame}

% Slide 16: Analisis Konvergensi Energi
\begin{frame}{Analisis Konvergensi Energi}
    \begin{itemize}
        \item Laju penurunan energi menunjukkan stabilitas algoritma SPSA.
        \item Perbandingan iterasi: GT-VQE mencapai konvergensi lebih cepat dibanding inisialisasi random.
    \end{itemize}
    % \includegraphics[width=0.6\textwidth]{Convergence_Graph.png}
\end{frame}

% Slide 17: Dinamika Internal (Entropy Analysis)
\begin{frame}{Dinamika Internal (Entropy Analysis)}
    \begin{itemize}
        \item \textbf{Entropi Von Neumann:} Mengukur derajat \textit{entanglement} and kemurnian state.
        \item \textbf{Fenomena Collapsing:} Penurunan entropi menuju nol mengindikasikan transisi state menuju probabilitas tunggal (solusi optimal).
    \end{itemize}
\end{frame}

% Slide 16: Hasil Backtesting Portofolio (N=2  dan   N=4)
\begin{frame}{Hasil Backtesting Portofolio (N=2  dan   N=4)}
    \begin{table}[]
        \centering
        \begin{tabular}{@{}lccc@{}}
        \toprule
        Metrik dan  Strategi Klasik dan  VQE Murni dan  \textbf{GT-VQE (Proposed)} \\ \midrule
        Sharpe Ratio (N=4) dan  0.85 dan  0.92 dan  \textbf{1.0107} \\
        Cumulative Return dan  +12.4\% dan  +14.8\% dan  \textbf{+18.5\%} \\
        Max Drawdown dan  -5.2\% dan  -4.1\% dan  \textbf{-3.8\%} \\ \bottomrule
        \end{tabular}
        \caption{Perbandingan performa portofolio selama periode 2021-2023.}
    \end{table}
\end{frame}

% Slide 19: Analisis Jendela Waktu Ekstrem
\begin{frame}{Analisis Jendela Waktu Ekstrem}
    \begin{itemize}
        \item Algoritma diuji pada periode volatilitas tinggi.
        \item \textbf{Agilitas Rebalancing:} Mekanisme GT-VQE mampu menyesuaikan bobot aset secara responsif terhadap guncangan pasar.
    \end{itemize}
\end{frame}

% Slide 20: Perbandingan Efisiensi Kedalaman Sirkuit
\begin{frame}{Perbandingan Efisiensi Kedalaman Sirkuit}
    \begin{itemize}
        \item Kedalaman optimal ditemukan pada $L=2$ hingga $L=3$.
        \item Penambahan layer lebih lanjut tidak memberikan peningkatan akurasi signifikan namun menambah akumulasi noise.
    \end{itemize}
\end{frame}
\section{Penutup}

% Slide 21: Kesimpulan
\begin{frame}{Kesimpulan}
    \begin{enumerate}
        \item \textbf{Formulasi EPG-Hamiltonian:} Berhasil memetakan utilitas ekonomi ke dalam sistem fisis secara konsisten.
        \item \textbf{Strategi Warm-Start:} Penggunaan PSNE terbukti efektif memandu konvergensi and mengatasi \textit{Barren Plateaus} pada sirkuit kuantum.
        \item \textbf{Superioritas Performa:} Model GT-VQE menghasilkan performa finansial yang lebih tangguh dengan \textit{Sharpe Ratio} 1,0107.
    \end{enumerate}
\end{frame}

% Slide 22: Saran  dan   Penutup
\begin{frame}{Saran  dan   Penutup}
    \begin{block}{Saran Pengembangan}
        \begin{itemize}
            \item Pengujian pada infrastruktur perangkat keras kuantum riil (bukan simulator).
            \item Integrasi model noise (\textit{decoherence, gate fidelity}) yang lebih realistis.
            \item Perluasan aset pada pasar yang mengalami fenomena \textit{Black Swan}.
        \end{itemize}
    \end{block}
    \vspace{1cm}
    \begin{center}
        \huge \textbf{Terima Kasih} \\
        \medskip
        \large Ada Pertanyaan?
    \end{center}
\end{frame}

\end{document}
